% =====================================================================
%  2bat-ud6-3.tex  —  SESSIÓ 3: Biaixos de mostreig: quan una mostra enganya
%  (sense preàmbul: s'inclou des de main.tex)
% =====================================================================
\horatitol{Sessió 3 — Biaixos de mostreig: quan una mostra enganya}%
{Aules pensants: davant d'una col·lecció de casos, detectar per què cada mostra pot enganyar i en quina direcció esbiaixaria el resultat.}

\subsection{Idea clau}

Aprofundim, a les \textbf{pissarres verticals} i en grups, en una de les claus de la
unitat: \textbf{per què una mostra mal triada porta a conclusions falses}. Per a cada cas
haureu de detectar \emph{per què} la mostra enganya, \emph{en quina direcció} (el resultat
sortiria més alt o més baix del real) i \emph{com} es podria arreglar.

\begin{idea}
\textbf{Tres tipus de biaix de mostreig habituals}
\begin{itemize}[leftmargin=1.4em, itemsep=3pt]
  \item \textbf{Biaix de selecció:} la manera de triar afavoreix un perfil concret (per
        exemple, enquestar només qui passa per un lloc determinat).
  \item \textbf{Biaix de no-resposta:} responen sobretot les persones amb una opinió molt
        marcada (a favor o en contra), i la resta callen.
  \item \textbf{Biaix de cobertura:} el mètode no arriba a tota la població (per exemple,
        una enquesta en línia no arriba a qui no té internet).
\end{itemize}
\textbf{Per a cada cas:} qui queda fora? El resultat sortiria \emph{més alt} o \emph{més
baix} del real? Com es podria corregir?
\end{idea}

\subsection{Exemples}

\begin{exemple}[biaix de no-resposta]
Un servei envia una enquesta de satisfacció per correu i només la retornen un $10\%$ dels
usuaris. Sovint responen sobretot els \textbf{molt satisfets} (per agrair) o els
\textbf{molt insatisfets} (per queixar-se), mentre que la majoria «tèbia» no contesta. El
resultat pot quedar \textbf{polaritzat} i no reflectir l'opinió mitjana real: és un
\textbf{biaix de no-resposta}.
\end{exemple}

\begin{exemple}[generalitzar des d'un únic lloc]
Un estudi d'opinió es fa \textbf{en un únic barri} benestant i se'n generalitza el resultat
a \textbf{tota la ciutat}. Com que aquell barri no representa la diversitat de la ciutat, la
mostra té \textbf{biaix de selecció}: el resultat descriuria aquell barri, no la ciutat.
Per corregir-ho, caldria mostrejar barris diversos en la proporció adequada.
\end{exemple}

\subsection{Exercicis resolts}

\begin{exresolt}[direcció del biaix]
Per mesurar la satisfacció amb un curs, s'enquesta només qui l'ha \textbf{acabat}. En quina
direcció s'esbiaixa el resultat i per què?
\solucio
Qui va abandonar el curs —potser per insatisfacció— no és a la mostra. Així, s'enquesta
sobretot qui hi estava prou a gust per acabar-lo. El resultat \textbf{sobreestima} la
satisfacció (surt \emph{més alt} del real). Per corregir-ho, caldria incloure també qui va
abandonar.
\resposta{Sobreestima la satisfacció, perquè exclou qui va abandonar.}
\end{exresolt}

\begin{exresolt}[identificar el tipus de biaix]
Una enquesta telefònica es fa només a \textbf{telèfons fixos}, en horari laboral. Quin
tipus de biaix hi ha i qui en queda fora?
\solucio
És un \textbf{biaix de cobertura} (i de selecció): els telèfons fixos en horari laboral
deixen fora qui no en té (sovint gent jove) i qui treballa i no pot respondre. La mostra
sobrerepresenta un perfil concret (per exemple, persones grans o a casa en aquell horari),
de manera que el resultat no és generalitzable a tota la població.
\resposta{Biaix de cobertura: en queden fora qui no té fix i qui treballa en aquell horari.}
\end{exresolt}

\subsection{Exercicis per fer}

\noindent\textit{Treballeu en grup a la pissarra: per a cada cas, digueu qui queda fora, en
quina direcció esbiaixa i com s'arreglaria.}

\begin{exercicis}
\begin{exlist}
  \item Una web de notícies fa una enquesta clicable («Estàs a favor de X?») i publica el
        resultat com l'opinió de la població. Quin biaix hi ha? Qui respon aquestes
        enquestes?
  \item Per saber quant llegeixen els joves, s'enquesta a la sortida d'una biblioteca. En
        quina direcció s'esbiaixaria el resultat?
  \item Un estudi sobre l'ús del cotxe enquesta persones en una àrea de servei d'autopista.
        Qui queda fora? Sobreestima o infraestima l'ús del cotxe?
  \item Per a cada situació, digues quin tipus de biaix predomina (selecció, no-resposta o
        cobertura):
        \begin{enumerate}[label=\alph*), leftmargin=1.6em, topsep=2pt]
          \item enquesta en línia sobre necessitats de la gent gran;
          \item enquesta on només responen els molt descontents.
        \end{enumerate}
  \item \textbf{(Correcció)} Tria un dels casos anteriors i explica com el faries més
        representatiu.
  \item \textbf{(Raonament)} Per què una mostra \textbf{gran} però \textbf{esbiaixada} no
        és millor que una de petita però ben triada? Pensa en un exemple.
\end{exlist}
\end{exercicis}

% ---------------------------------------------------------------------
\fullsolucions{Sessió 3}
\begin{solucions}
\begin{sollist}
  \item Biaix de selecció i de no-resposta: responen només els lectors de la web
        interessats en el tema i motivats a clicar (autoseleccionats). No és una mostra
        representativa de la població.
  \item Sobreestimaria la lectura: a la sortida d'una biblioteca gairebé tothom llegeix
        (mostra esbiaixada cap a lectors).
  \item En queda fora qui no viatja en cotxe (o no per autopista). Sobreestima l'ús del
        cotxe, perquè s'enquesta justament qui l'està fent servir.
  \item (a) biaix de cobertura (no arriba a qui no té internet); (b) biaix de no-resposta
        (responen sobretot els descontents).
  \item Resposta oberta. Ha de proposar ampliar o diversificar la manera d'arribar a la
        població (diferents llocs, moments, canals) i incloure els grups que quedaven fora.
  \item Resposta oberta. Idea: una mostra gran però esbiaixada repeteix el mateix error a
        gran escala (mesura molt bé un grup equivocat); la mida no corregeix el biaix.
        Exemple: un milió de respostes d'una enquesta clicable segueix representant només
        qui clica, no la població.
\end{sollist}
\end{solucions}
